\documentclass[11pt]{article}
\usepackage[margin=1in]{geometry}
\usepackage{amsmath,amssymb,amsthm,hyperref}
\setlength{\emergencystretch}{2em}
\newtheorem{theorem}{Theorem}
\newtheorem{lemma}[theorem]{Lemma}
\newcommand{\R}{\mathbb R}
\title{Minimum means and the literal both-types invariance classification}
\author{Myeongjin Kang}
\date{October 26, 2026}
\begin{document}
\maketitle
\begin{abstract}
The binary and ternary minimum functions are continuous symmetric means on the positive orthant and satisfy both invariance identities of Horwitz. They cannot be simultaneous inverse transforms of the binary and ternary arithmetic means under any genuine invertible generator. This gives a complete counterexample to the literal unrestricted continuous-mean formulation recorded as MathDB 376980. The qualification is essential: the witness is not strict, and the result does not address a classification with an added strictness or analyticity assumption. A single Lean module verifies the source mean definitions, all-permutation symmetry, both invariance identities, and the inverse-generator obstruction.
\end{abstract}
\noindent\textit{Credit and license.} This AI-assisted note verifies an elementary minimum construction without a novelty claim. Original prose is CC BY 4.0 and original code is MIT.


\section{The exact source statement and its scope}
In Section 1 of \emph{Invariant Means}, Horwitz defines a mean on $\R_{+}^{n}$ to be continuous and to lie between the minimum and maximum of its arguments. Strictness is then defined as a separate property: an endpoint is attained only if all arguments agree \cite{horwitz}. The source also defines symmetry under every permutation.

For a binary mean $m$ and ternary mean $M$, type 1 invariance is
\begin{equation}\label{eq:type1}
 M(m(a,c),m(a,b),m(b,c))=M(a,b,c),\qquad a,b,c>0,
\end{equation}
and type 2 invariance is
\begin{equation}\label{eq:type2}
 M(a,b,m(a,b))=m(a,b),\qquad a,b>0.
\end{equation}
In Section 7.4 the original PDF's Conjecture 29, on printed page 26, asserts that the only pairs of means satisfying both types are
\begin{align}
 \overline m(a,b)&=h^{-1}\left(\frac{h(a)+h(b)}2\right),\label{eq:binary}\\
 \overline M(a,b,c)&=h^{-1}\left(\frac{h(a)+h(b)+h(c)}3\right),\label{eq:ternary}
\end{align}
for a monotonic function $h$ on $(0,\infty)$. The arXiv HTML conversion renumbers this statement as Conjecture 26; the displayed wording agrees. MathDB 376980 records this classification \cite{mathdb}.

The literal conjecture does not add strictness or analyticity to the Section 1 definition of a mean. Earlier results impose such extra properties explicitly; in particular, the strictness assumption in Section 3 is stated for that section. Our conclusion concerns the unrestricted formulation just described. It does not refute an intended stronger statement confined to strict means, strictly isotone means, or analytic means.

The argument below uses familiar elementary properties of the minimum. It is supplied as a verified source-level counterexample, not as a new-to-literature theorem or a claim of priority. The archived MathDB record has no posted solution to attribute; this absence is not evidence of novelty.

\section{A complete counterexample}
Define, on all real inputs and hence on the positive orthants,
\[
 m_{\min}(a,b)=\min(a,b),\qquad
 M_{\min}(a,b,c)=\min(a,b,c)=\min(a,\min(b,c)).
\]
\begin{theorem}\label{thm:main}
The pair $(m_{\min},M_{\min})$ consists of continuous symmetric means satisfying both \eqref{eq:type1} and \eqref{eq:type2}. It admits no representation \eqref{eq:binary}--\eqref{eq:ternary} with a genuine inverse $h^{-1}$, even if monotonicity of $h$ is dropped. Thus it refutes the literal classification of all continuous symmetric means.
\end{theorem}
\begin{proof}
The function $(a,b)\mapsto\min(a,b)$ is continuous on $\R^2$. Composing two such functions proves continuity of $M_{\min}$ on $\R^3$. These global continuity assertions imply continuity on the required positive orthants.

The internality bounds are immediate:
\[
 \min(a,b)\le m_{\min}(a,b)\le\max(a,b),
\]
\[
 \min(a,b,c)\le M_{\min}(a,b,c)\le\max(a,b,c).
\]
The lower bounds are equalities. Their values are positive when all inputs are positive, so the substituted arguments in both invariance equations stay in the source domain. The minimum of a finite list is unchanged by any permutation, proving symmetry in both arities.

For type 1 invariance, evaluate the actual source argument order:
\begin{align*}
&M_{\min}(m_{\min}(a,c),m_{\min}(a,b),m_{\min}(b,c))\\
&\hspace{1cm}=\min(\min(a,c),\min(a,b),\min(b,c))
 =\min(a,b,c)=M_{\min}(a,b,c).
\end{align*}
The equality holds because each of $a,b,c$ occurs in the nested minimum, with repetitions irrelevant. For type 2,
\[
 M_{\min}(a,b,m_{\min}(a,b))=\min(a,b,\min(a,b))=\min(a,b)=m_{\min}(a,b).
\]
Both identities hold for all real inputs, in particular on the positive domains.

It remains to rule out the actual inverse-transform conclusion, rather than just a selected proposed generator. Any genuine inverse representation \eqref{eq:binary} has an injective $h$ on $(0,\infty)$, and applying $h$ gives
\begin{equation}\label{eq:generator}
 h(m_{\min}(a,b))=\frac{h(a)+h(b)}2\qquad(a,b>0).
\end{equation}
At $a=1,b=2$, this says
\[
 h(1)=\frac{h(1)+h(2)}2,
\]
whence $h(1)=h(2)$, contradicting injectivity. Thus even the binary formula is impossible, so certainly no simultaneous pair representation exists. Monotonicity was not needed in the contradiction.
\end{proof}

\section{The inverse-formula bridge in detail}
For clarity about the meaning of $h^{-1}$, the formal proof does not silently treat an arbitrary function as an inverse. It uses a class of representations larger than the one in the source. Suppose $h,\iota\colon\R\to\R$ satisfy
\begin{equation}\label{eq:leftinverse}
 \iota(h(a))=a\qquad(a>0).
\end{equation}
Require also that each binary and ternary arithmetic average in \eqref{eq:binary} and \eqref{eq:ternary} belongs to $h((0,\infty))$. This is exactly the domain requirement for evaluating an ordinary inverse on those averages. Finally, require the two displayed representation formulas with $\iota$ in place of $h^{-1}$. We impose no monotonicity, continuity, or global surjectivity on $h$.

Equation \eqref{eq:leftinverse} already implies injectivity of $h$ on the positive domain: applying $\iota$ to $h(a)=h(b)$ gives $a=b$. If $u\in h((0,\infty))$, write $u=h(c)$ with $c>0$; then
\[
 h(\iota(u))=h(\iota(h(c)))=h(c)=u.
\]
Consequently, applying $h$ to the inverse formulas yields \eqref{eq:generator} and its correctly normalized ternary counterpart,
\[
 h(M(a,b,c))=\frac{h(a)+h(b)+h(c)}3.
\]
The formal proof derives both equations from the inverse representation and its range conditions. No missing inverse, injectivity, or arithmetic-conjugacy bridge is assumed. Since the minimum pair fails even this enlarged class, it also fails every monotone genuinely invertible instance of the source conclusion.

\section{Why the qualification matters}
The witness is an allowed mean in the source's basic sense, but it is not strict:
\[
 m_{\min}(1,2)=1,\qquad 1\ne2.
\]
Its lack of strictness is the reason it evades ordinary quasi-arithmetic classification. This is an explicit mathematical obstruction, not a counterexample to a stricter statement that includes this extra hypothesis. The source's motivation uses differentiable logarithmic means and derivative calculations; our theorem does not challenge those calculations or resolve a conjecture restricted to their smooth or strict setting.

The only source assertion decided here is the written classification of all continuous symmetric means satisfying the two identities. One cannot infer from this argument either that a strict classification is true or that it is false. The counterexample also does not claim anything about whether every given mean has a type 1 or type 2 invariant extension.

\section{Formal coverage and verification}
The file \texttt{formal/InvariantMinimum.lean}, namespace \texttt{InvariantMinimum}, contains the entire proof. Its \texttt{IsMean2} and \texttt{IsMean3} predicates express continuity on the positive orthants and the source minimum/maximum bounds. Its symmetry predicates quantify over every permutation of \texttt{Fin 2} and \texttt{Fin 3}. The definitions \texttt{Type1} and \texttt{Type2} use precisely \eqref{eq:type1} and \eqref{eq:type2}. The witness functions are explicit nested minima; no source-class membership or invariance identity is assumed.

The predicate \texttt{InverseArithmeticPair} spells out inverse cancellation on positive inputs, membership of the two averages in the generator's range, and both inverse formulas. The theorem \texttt{inverse\_pair\_implies\_generator} proves the conjugacy bridge. The theorem \texttt{no\_arithmetic\_generator} supplies the $(1,2)$ obstruction, and \texttt{no\_inverse\_arithmetic\_pair} combines the two. The main theorem \texttt{source\_counterexample} bundles continuous mean membership, full symmetry, both invariance identities, and the exclusion of the inverse-transform conclusion, with no unproved bridge hypotheses.

The pinned campaign verifier compiled the source and a separate all-declaration audit under Lean 4.34.1 with warnings treated as errors. All 24 declarations passed. The observed foundational dependencies are exactly \texttt{propext}, \texttt{Classical.choice}, and \texttt{Quot.sound}. The exact source hash and dependency validation receipt are in \texttt{build/audit.json}; source snapshots and statement-label discrepancies are recorded in the provenance files. The accompanying independent review checks both formal coverage and source scope.

\begin{thebibliography}{9}
\bibitem{horwitz} Alan Horwitz, \emph{Invariant Means}, arXiv:math/0007095v1, July 15, 2000. Section 1, printed pp.~1--2; Section 7.4, Conjecture 29, printed p.~26 of the original PDF. The HTML conversion labels the conjecture 26. \url{https://arxiv.org/abs/math/0007095v1}.
\bibitem{mathdb} MathDB 376980, \emph{Characterization of pairs invariant in both types}, \url{https://mathdb.com/p/376980}. The exact post identifier and source hashes accompany the certificate.
\end{thebibliography}
\end{document}
